\hypersetup{pageanchor=false}
\begin{titlepage}
\centering
{\large University of Sheffield\par}
\vspace{0.25cm}
{\normalsize Department of Computer Science\par}
\vspace{3cm}
{\LARGE\bfseries VolPFN: Volatility Forecasting\par}
\vspace{0.35cm}
{\LARGE\bfseries from a Synthetic Prior\par}
\vspace{1cm}
{\large Performance, prior attribution and tail recalibration\par}
\vspace{2.8cm}
{\Large Bader Matar\par}
\vspace{0.7cm}
{\normalsize MSc in Artificial Intelligence\par}
\vspace{1.1cm}
{\normalsize Supervisor: Georgios Moulantzikos\par}
\vspace{0.7cm}
{\normalsize Registration number: \rule{3cm}{0.4pt}\par}
\vfill
{\normalsize Dissertation submitted for the MSc in Artificial Intelligence\par}
\vspace{0.3cm}
{\normalsize September 2026\par}
\vspace{0.3cm}
{\footnotesize Research revision: 2 October 2026; Stage 4 editorial cleanup: 3 October 2026\par}
\end{titlepage}
\hypersetup{pageanchor=true}

\chapter*{Declaration}
\addcontentsline{toc}{chapter}{Declaration}
All sentences or passages quoted in this report from other people's work have been specifically acknowledged by clear cross-referencing to author, work and page(s). Any illustrations that are not the work of the author of this report have been used with the explicit permission of the originator and are specifically acknowledged. I understand that failure to do this amounts to plagiarism and will be considered grounds for failure in this project and the degree examination as a whole.

\vspace{1.5cm}
\noindent Name: Bader Matar

\vspace{0.8cm}
\noindent Signature:

\vspace{0.8cm}
\noindent Date:



\chapter*{Abstract}
\addcontentsline{toc}{chapter}{Abstract}
Financial volatility changes over time, while short market histories provide limited examples for learning. We develop VolPFN, a neural forecaster trained on simulated returns, to examine what synthetic training contributes to market prediction. Its weights remain fixed during use; calibration and forecast combinations are fitted separately.

The main experiment matches architecture, initial weights and training allocation across simulation designs. Broad training improves the common log score by about 0.0052 nats per return, with an interval supporting a positive gain. The improvement is modest and concentrated in a few markets. It supports the combined design without identifying an individual causal ingredient.

VolPFN leads the tested general pretrained benchmark, and its variants lead the daily-return-only leaderboard. These rankings evaluate complete procedures without establishing adjusted superiority over every competitor. Distribution prediction outperforms direct risk prediction, while extra volatility supervision has no detected benefit. Regularisation and development selection make historical training competitive. A daily--intraday blend improves on calibrating the intraday model alone, with stronger alternatives unresolved. On fixed exchange-traded funds, synthetic training without extra supervision improves on the historical procedure, but superiority over the rolling financial benchmark remains inconclusive.

A compact fitted rule retains much of an earlier network's scale advantage without explaining its computation. Residual correction helps most clearly with Gaussian tails. Additional neural joint-score gains remain uncertain, and the fixed correction rule sometimes loses to filtered historical simulation.

The evidence therefore supports useful prediction within defined comparisons. Prior data exposure, shared shocks and uncertainty conditional on saved forecasts limit generalisation.

\chapter*{Acknowledgements}
\addcontentsline{toc}{chapter}{Acknowledgements}
I thank Georgios Moulantzikos for his supervision and the Kuwait government scholarship programme for supporting my studies. I also thank the researchers who provided the Oxford-Man data and the replication material used in this project.

\chapter*{Abbreviations and reading guide}
\addcontentsline{toc}{chapter}{Abbreviations and reading guide}
The reading guide assumes familiarity with averages and graphs without prior knowledge of financial volatility or prior-fitted networks. Chapter~\ref{ch:background} introduces the financial and machine-learning concepts needed for the experiments. The appendices retain the detailed implementation specifications and supporting results.

\begingroup
\small
\begin{longtable}{@{}>{\raggedright\arraybackslash}p{2.7cm}>{\raggedright\arraybackslash}p{11.0cm}@{}}
\toprule Term & Meaning in this dissertation \\
\midrule\endfirsthead
\toprule Term & Meaning in this dissertation (continued) \\
\midrule\endhead
ARCH / GARCH & Models that update a changing variance from past shocks and, for GARCH, earlier variance estimates. \\
EWMA & Exponentially weighted moving average: gives progressively less weight to older observations. \\
LSTM & Long short-term memory: a recurrent neural architecture used in the daily comparisons. \\
RK & Realised kernel: an estimate of variability within a trading session, used as a proxy or a past input. \\
HAR & Heterogeneous autoregressive model: combines daily, weekly and monthly realised-risk summaries. \\
STES & Smooth-transition exponential smoothing: changes the updating response with observed conditions. \\
GNNHAR & Graph neural HAR: combines realised-risk summaries using links between assets. \\
TTM & Tiny Time Mixers: the pretrained model supplying a forecast from past realised kernels. \\
GINN & GARCH-informed neural network: combines observed innovations and a fitted GARCH teacher. \\
GP & Gaussian process: a probability model over functions. \\
PFN / VolPFN & Prior-data fitted network; VolPFN is the volatility-oriented network developed here. \\
QLIKE & A loss for evaluating positive risk forecasts; smaller values indicate better prediction. \\
RMS & Root-mean-square scale used to normalise returns; the stated centring and observation rules depend on the study. \\
M2 & Expected squared return (the second moment). It equals variance plus squared predictive mean. \\
NLL / nats & Negative log likelihood, with lower values preferred; nats are log-score units using natural logarithms. \\
ETF & Exchange-traded fund: a fund traded on an exchange, used here to extend the evaluated asset population. \\
VaR / ES & Value-at-Risk is a lower-return threshold; Expected Shortfall is the average return below that threshold. \\
FZ0 & A loss that evaluates VaR and ES together; smaller is better. \\
FHS & Filtered historical simulation: estimate a residual distribution after removing scale, then restore forecast scale. \\
CI & Confidence interval: uncertainty under the stated sampling procedure and assumptions. \\
Arm & One experimental configuration; it specifies the model and training procedure being compared. \\
Full / noaux & Full uses extra supervision from simulated volatility labels. Noaux removes that training loss while retaining the same physical output heads. \\
Direct / historical & Direct predicts expected squared return with QLIKE training. Historical learns the return distribution from earlier market windows. \\
Fixed-final & Use the checkpoint saved at the final scheduled update, without choosing it from validation scores. \\
Checkpoint / seed & A checkpoint stores trained weights. A seed sets pseudorandom training choices. \\
TRAIN / development & TRAIN denotes the observations used for fitting. Development observations are used for the stated selection or calibration decisions. \\
Common-$t_8$ / native & Common-$t_8$ evaluates scales under the same Student distribution with eight degrees of freedom. Native scoring evaluates each model's own return distribution. \\
UniVol & The project's non-neural forecast that combines updating risk filters and restart experts. \\
Compact rule / distillation & A small explicit predictor; distillation fits that predictor to imitate network forecasts. Likelihood fitting instead predicts observed returns. \\
Helper / blend & A helper supplies an additional forecast. A blend combines it with the daily forecast using parameters fitted from earlier observations. \\
Scoring mask / key & The mask fixes eligible observations. A key identifies one market and date. \\
Union date / family & A union date occurs in at least one included market. A comparison family is the set of tests adjusted together for multiplicity. \\
Guard / hindsight & The guard follows the study's fixed rule for applying tail correction. Hindsight alternatives use later information and are diagnostic. \\
Oracle / reference & An oracle uses known simulated truth. Other references may also use information unavailable to the evaluated forecaster; their privileges are specified. \\
\bottomrule
\end{longtable}
\endgroup
